% Generated by scripts/export_feature_tables.py -- do not edit by hand.
\begin{description}
  \item[\texttt{user\_activity\_count}] Review volume separates a user with enough history for collaborative signal from one riding on the popularity prior. The distribution is heavy-tailed, so it is log-compressed before NMF's input matrix or any downstream linear model consumes it.

  \item[\texttt{user\_mean\_rating}] Implicit filtering keeps only positives, so this measures how generously a user rates recipes they already liked. It calibrates the confidence attached to a positive rather than the positive itself.

  \item[\texttt{user\_tenure\_days}] Calendar span between a user's first and last training review. Distinguishes fifty reviews spread over three years from fifty reviews in one weekend; the two imply very different taste stability.

  \item[\texttt{user\_recency\_days}] Days from a user's last training review to the cut-off. A dormant user's stored preferences are staler than an active user's, and the serving path should weigh them accordingly. Depends on the cut-off, so it is recomputed per partition and never cached across splits.

  \item[\texttt{item\_popularity\_rank}] Normalized dense popularity rank over the full recipe catalogue. Backs the popularity baseline and the cold-start fallback. Total over the catalogue by construction, so serving never hits a missing key.

  \item[\texttt{item\_trailing\_popularity}] Review count in the 90 days before the cut-off. Lifetime popularity is dominated by whatever recipe went viral years ago; a trailing window tracks what is popular now, which is what a system serving traffic at the cut-off actually needs. This is the feature that justifies evaluating on a temporal split.

  \item[\texttt{item\_category\_tokens}] Recipe category label (Dessert, Main Dish, Salad, ...). Content signal available for recipes with no interactions at all, which is the only thing a purely collaborative model cannot supply for a cold recipe.

  \item[\texttt{item\_calories}] Calories per serving, passed through from the recipe catalogue. Feeds the nutrition-aware post-filter directly rather than the NMF ranker, so the domain constraint stays independent of whatever the collaborative model has learned.

  \item[\texttt{item\_sodium\_mg}] Sodium content in milligrams, passed through from the catalogue. Together with calories it is the minimal nutrition signal the post-filter needs to enforce hard dietary constraints on the ranked candidate list.

\end{description}
